\documentclass[10pt]{article}

% ============================================================
% NeurIPS-style formatting (compact, single column)
% ============================================================
\usepackage[utf8]{inputenc}
\usepackage[T1]{fontenc}
\usepackage{times}
\usepackage[margin=0.75in,top=0.7in,bottom=0.7in]{geometry}
\usepackage{microtype}

% Math
\usepackage{amsmath,amssymb,amsthm}
\newtheorem{proposition}{Proposition}

% Tables and figures
\usepackage{booktabs}
\usepackage{tabularx}
\usepackage{multirow}
\usepackage{graphicx}
\usepackage{caption}
\captionsetup{font=small,skip=4pt}

% Lists
\usepackage{enumitem}
\setlist{nosep,leftmargin=*}

% Hyperlinks
\usepackage[colorlinks=true,citecolor=blue,linkcolor=blue,urlcolor=blue]{hyperref}

% Citations (compact)
\usepackage{natbib}
\bibliographystyle{abbrvnat}
\setlength{\bibsep}{1pt plus 0.5pt}

% Compact spacing
\usepackage{setspace}
\setstretch{0.94}
\usepackage{titlesec}
\titlespacing*{\section}{0pt}{6pt plus 2pt minus 2pt}{3pt plus 1pt minus 1pt}
\titlespacing*{\subsection}{0pt}{5pt plus 2pt minus 1pt}{2pt plus 1pt minus 1pt}
\titlespacing*{\paragraph}{0pt}{3pt plus 1pt minus 1pt}{0.5em}
\titleformat{\section}{\normalfont\large\bfseries}{\thesection}{0.5em}{}
\titleformat{\subsection}{\normalfont\normalsize\bfseries}{\thesubsection}{0.5em}{}

% Reduce spacing around equations
\AtBeginDocument{
  \abovedisplayskip=4pt plus 2pt minus 2pt
  \belowdisplayskip=4pt plus 2pt minus 2pt
  \abovedisplayshortskip=2pt plus 1pt
  \belowdisplayshortskip=2pt plus 1pt
}

% Table cell vertical padding
\renewcommand{\arraystretch}{1.05}

% Checkmark and cross symbols
\usepackage{pifont}
\newcommand{\cmark}{\ding{51}}
\newcommand{\xmark}{\ding{55}}

% Reduce float spacing
\setlength{\textfloatsep}{6pt plus 2pt minus 2pt}
\setlength{\floatsep}{4pt plus 2pt minus 2pt}
\setlength{\intextsep}{4pt plus 2pt minus 2pt}

% ============================================================
% Title
% ============================================================
\title{\Large\bfseries Retrieval-Augmented Input-Space Translation for Frozen Clinical Predictors}

\author{
  \textit{NeurIPS 2026 --- Positioning Document}
}

\date{}

\begin{document}

\maketitle
\vspace{-10pt}

\begin{abstract}
\noindent\small
Clinical prediction models validated at one hospital degrade when deployed at another due to distributional shift in electronic health records (EHR).
Standard domain adaptation methods require modifying the target model---incompatible with regulatory frameworks that mandate locked clinical algorithms.
We propose \emph{input-space translation} for frozen clinical predictors: learning a neural translator that maps source EHR time series into the target input distribution while the predictor remains \emph{entirely untouched}.
We develop three architectures of increasing capability---delta-based residual, shared latent encoding, and retrieval-augmented cross-attention grounded in a target-domain memory bank---and introduce the \emph{gradient alignment condition} $\alpha(\theta) = \cos(\nabla\mathcal{L}_{\text{task}}, \nabla\mathcal{L}_{\text{fid}})$, a diagnostic measurable in the first training epoch that predicts when frozen-model translation will succeed or fail.
Across five clinical tasks on eICU$\to$MIMIC-IV, the retrieval translator achieves universal improvements of up to $+5.6$ AUROC and $+16.1$ AUCPR over the frozen baseline: on AKI, it surpasses even the \emph{in-domain} MIMIC-native LSTM in both AUROC (91.14 vs.\ 89.7) and AUCPR (72.9 vs.\ 66.5), all three classification tasks exceed the eICU-native LSTM on both metrics, and length-of-stay regression matches the eICU-native model---all with zero modification to the frozen predictor.
To the best of our knowledge, this is the first work to achieve cross-hospital domain adaptation for clinical time series while keeping the target predictor entirely frozen, enabled by instance-level retrieval from the target domain and a gradient alignment theory that predicts when frozen-model translation succeeds.
\end{abstract}

\vspace{-6pt}

% ============================================================
% Part 1: Related Work
% ============================================================
\section{Related Work and Current Methods}
\label{sec:related}

\paragraph{Domain adaptation in clinical and general time series.}
Domain adaptation for electronic health records has been approached through
representation-space alignment~\cite{purushotham2017vrada,zhang2022adadiag,extracare2026},
optimal transport~\cite{ottehr2025},
inductive transfer with model fine-tuning~\cite{mutnuri2024,driever2025},
and data harmonization pipelines~\cite{purushotham2018benchmark,harutyunyan2019multitask}.
VRADA~\cite{purushotham2017vrada} introduced variational RNNs with adversarial alignment for
clinical time series but requires end-to-end training of the target model.
OTTEHR~\cite{ottehr2025} applies optimal transport to the same eICU--MIMIC databases we study,
yet performs population-level feature alignment without temporal modeling and without any
frozen-model constraint.
ExtraCare~\cite{extracare2026}, a concurrent work on interpretable clinical DA through concept-grounded
orthogonal decomposition, also operates in representation space.
In the broader time-series DA literature, established methods---CoDATS~\cite{wilson2020codats},
RAINCOAT~\cite{he2023raincoat}, CLUDA~\cite{ozyurt2023cluda},
ACON~\cite{liu2024acon}---learn domain-invariant representations through a shared feature
extractor, typically with adversarial~\cite{wilson2020codats}, contrastive~\cite{ozyurt2023cluda},
or spectral kernel~\cite{liu2021advskm} objectives.
ADATIME~\cite{ragab2023adatime} benchmarks eleven such methods across fifty scenarios, all
operating in representation space.
MAPU~\cite{ragab2023mapu} addresses the orthogonal constraint of restricted \emph{data} access
(source-free adaptation) but still modifies the target model.
A unifying limitation of every method above is the assumption that the target model's
parameters are trainable: they either fine-tune the predictor or jointly learn a feature
extractor.
When the target predictor must remain frozen---as required for regulatory compliance with
locked clinical algorithms under the FDA's Software as a Medical Device
(SaMD) framework~\cite{fda2021samd}---these methods are inapplicable,
motivating input-space adaptation approaches.

\paragraph{Frozen-model adaptation and input-space methods.}
A growing body of work modifies \emph{inputs} rather than model weights.
Prompt tuning~\cite{lester2021prompt} learns a fixed soft-token prefix for frozen language models;
Visual Prompt Tuning (VPT)~\cite{jia2022vpt} extends this to frozen vision transformers.
Adversarial reprogramming~\cite{elsayed2019reprogramming} repurposes frozen ImageNet classifiers
via a single universal additive perturbation.
Voice2Series~\cite{yang2021voice2series} reprograms frozen acoustic models for time-series
classification, and Time-LLM~\cite{jin2024timellm} reprograms frozen LLMs (GPT-2, LLaMA) for
time-series forecasting---the most prominent frozen-model work in the time-series literature,
with over 1000 citations.
TATO~\cite{qiu2026tato}, a concurrent ICLR~2026 contribution, explicitly proposes
``adapt data to model'' for frozen time-series foundation models using handcrafted
transformation operators (slicing, normalization, outlier correction).
L2C~\cite{chi2025l2c} learns input-space modifications for frozen CLIP at test time.
PixelDA~\cite{bousmalis2017pixelda} performs input-space DA in vision via GANs but does
not freeze the downstream model.
SHOT~\cite{liang2020shot} partially freezes the classifier head while adapting the feature
extractor; our constraint is stronger in that the \emph{entire} model is frozen.
Table~\ref{tab:frozen_comparison} summarizes the key distinctions.
These methods face three limitations in our setting.
First, prompt tuning, VPT, and reprogramming apply \emph{input-independent} transformations: they learn a single perturbation that is shared identically across all inputs.
While our translator also has fixed (learned) parameters, its output is \emph{input-conditioned}---each patient's time series receives a distinct transformation adapted to its specific values, length (1 to 169 timesteps), and missingness pattern across 48+ features.
A universal additive perturbation cannot provide such sample-level adaptation, as the correction needed for one patient's lab values may be entirely inappropriate for another's.
Second, the vast majority target \emph{task} adaptation (repurposing a model for a new task),
not \emph{domain} adaptation (translating data from one hospital's distribution to another's).
Third, handcrafted transforms (TATO) rely on a predefined set of operators (slicing, normalization, outlier correction) applied independently per feature.
Clinical domain shift, however, involves nonlinear, \emph{context-dependent} distributional differences: the correction required for a lab value depends on concurrent vitals, medications, and the patient's trajectory---interactions that a fixed pipeline of feature-wise operators cannot model.
To our knowledge, no prior frozen-model work simultaneously incorporates instance-level retrieval from the target domain,
evaluates on clinical EHR data, and provides theoretical grounding via gradient dynamics.
These capacity limitations motivate our
neural translator, which uses instance-level retrieval from the target domain for contextual
grounding.

\begin{table}[t]
\centering
\caption{Comparison of frozen-model and input-space adaptation methods. Our approach is the
only one that combines sample-conditional neural translation, target-domain retrieval,
variable-length support, and cross-domain DA for clinical time series.}
\label{tab:frozen_comparison}
\small
\setlength{\tabcolsep}{4pt}
\begin{tabular}{lccccc}
\toprule
\textbf{Method} & \textbf{Input modification} & \textbf{Sample-cond.} & \textbf{Retrieval} & \textbf{Var.-length} & \textbf{Purpose} \\
\midrule
Prompt Tuning~\cite{lester2021prompt} & Fixed prefix & \xmark & \xmark & \xmark & Task adapt. \\
VPT~\cite{jia2022vpt} & Prepended tokens & \xmark & \xmark & \xmark & Task adapt. \\
Voice2Series~\cite{yang2021voice2series} & Learned transform & \xmark & \xmark & \xmark & Cross-task \\
Time-LLM~\cite{jin2024timellm} & Reprogramming layer & \xmark & \xmark & Fixed patches & Cross-modality \\
TATO~\cite{qiu2026tato} & Handcrafted pipeline & Partially & \xmark & \cmark & Domain adapt. \\
\textbf{Ours} & \textbf{Neural translator} & \cmark & \cmark & \cmark & \textbf{Cross-domain DA} \\
\bottomrule
\end{tabular}
\end{table}

\paragraph{Retrieval-augmented models for time series.}
Retrieval-augmented generation (RAG)~\cite{lewis2020rag} established the paradigm of
conditioning neural models on retrieved external knowledge.
$k$NN-LM~\cite{khandelwal2020knnlm} augments a frozen language model with nearest-neighbor
retrieval from a token-level datastore, achieving domain adaptation by simply swapping the
datastore---without any model retraining.
$k$NN-MT~\cite{khandelwal2021knnmt} extends this to machine translation, performing
per-timestep $k$-NN lookup from a domain-specific datastore and interpolating
\emph{output token distributions} with the frozen model's predictions.
RETRO~\cite{borgeaud2022retro} introduced chunked cross-attention over retrieved document
passages---the closest architectural precedent for our fusion mechanism.
However, RETRO retrieves passages from the \emph{same} domain to augment language generation,
whereas we retrieve from the \emph{target} domain to transform \emph{source} inputs,
making retrieval an instrument of domain translation rather than knowledge augmentation.
In the time-series literature, retrieval augmentation is an emerging paradigm at top venues:
RATD~\cite{liu2024ratd} retrieves relevant series to guide diffusion-based forecasting
(NeurIPS~2024),
TS-RAG~\cite{ning2025tsrag} augments time-series foundation models with an adaptive
retrieval mixer (NeurIPS~2025),
and RAFT~\cite{raft2025} exploits inductive biases from retrieved examples for forecasting
(ICML~2025).
RAM-EHR~\cite{xu2024ramehr} retrieves textual medical knowledge from PubMed and
knowledge graphs to augment EHR predictions, but retrieves \emph{text}, not time-series data.
Our retrieval translator differs from all prior work in three respects:
(a)~we retrieve encoded \emph{target-domain time-series windows}---representations of what the
frozen predictor expects to see---not text passages or output distributions;
(b)~retrieved windows serve as \emph{cross-attention context} in a Transformer decoder,
enabling selective feature-level borrowing, rather than being interpolated with output
distributions ($k$NN-MT) or concatenated with the query (RAG); and
(c)~the purpose is cross-domain translation, not knowledge augmentation or forecasting.
$k$NN-MT~\cite{khandelwal2021knnmt} is the closest ancestor: both systems pair a frozen
model with per-step $k$-NN from a domain-specific store.
The critical difference is that $k$NN-MT interpolates output token distributions at
the decoder, operating within a single language, while our retrieval translator uses
retrieved embeddings as cross-attention context to transform the \emph{input}, bridging
fundamentally different data distributions (eICU vs.\ MIMIC-IV feature spaces with
heterogeneous recording practices, device calibrations, and documentation conventions).

\paragraph{Foundation models as an alternative paradigm.}
Large-scale pre-training offers a complementary path to cross-hospital generalization.
ICareFM~\cite{burger2025icarefm} pre-trains on 650K ICU stays from multiple hospitals with
a self-supervised time-to-event objective, achieving zero-shot transfer.
CLMBR/EHRSHOT~\cite{wornow2023ehrshot} provides a 141M-parameter EHR foundation model
with strong few-shot evaluation on 6,739 patients.
LLEMR~\cite{wu2024llemr} instruction-tunes an LLM on 400K+ MIMIC-IV examples.
Med-BERT~\cite{rasmy2021medbert} and BEHRT~\cite{li2020behrt} pre-train Transformers on
28.5M and 1.6M patient records, respectively.
We build directly on the YAIB benchmark framework~\cite{vandewater2024yaib}, whose
multi-center preprocessing pipeline and five standardized task definitions provide our
experimental infrastructure and frozen baselines (LSTMs in our experiments, though the input-space translation paradigm is architecture-agnostic).
Foundation models require (a)~massive multi-center datasets (100K--650K stays),
(b)~weeks of compute for pre-training, and (c)~replacing the deployed predictor
entirely, necessitating full regulatory re-validation.
Our approach is complementary, not competitive: when a hospital has a validated, deployed
predictor that cannot be retrained due to regulatory, institutional, or practical constraints,
input-space translation is the only viable adaptation strategy.
Foundation models and our translators address different points on the
\emph{deployment constraint spectrum}---from unconstrained model replacement to
fully frozen predictor preservation.

\paragraph{Domain adaptation theory and gradient dynamics.}
Ben-David et al.~\cite{bendavid2010} bound target-domain error by
$\varepsilon_T \leq \varepsilon_S + d_{\mathcal{H}\Delta\mathcal{H}}(\mathcal{S}, \mathcal{T}) + \lambda^*$,
where $d_{\mathcal{H}\Delta\mathcal{H}}$ measures distribution divergence and $\lambda^*$
is the ideal joint error.
Mansour et al.~\cite{mansour2009} refine this via discrepancy distance tailored to the
loss and hypothesis class; Zhang et al.~\cite{zhang2019mdd} bridge theory and algorithms
through margin disparity discrepancy.
Canonical methods operationalize these bounds:
DANN~\cite{ganin2016dann} minimizes $\mathcal{H}$-divergence via gradient reversal,
Deep CORAL~\cite{sun2016coral} aligns second-order feature statistics,
and DAN~\cite{long2015dan} matches multi-kernel MMD~\cite{gretton2012mmd} across layers.
All require a \emph{learnable} feature extractor.
In our frozen-model setting, representation-space alignment is impossible by construction: since $\nabla_\phi f_T = 0$, no gradient signal can reshape the frozen model's internal representations, and any alignment objective (MMD, adversarial, CORAL) applied to intermediate layers has no learnable pathway to reduce the domain gap.
We instead adapt these methods to operate in input space and evaluate them as baselines (detailed in the full paper).
%
The multi-objective structure of our training---minimizing task loss $\mathcal{L}_\text{task}$
(prediction accuracy through the frozen model) and fidelity loss $\mathcal{L}_\text{fid}$
(proximity to the source distribution) simultaneously---connects to the gradient conflict
literature in multi-task learning.
Sener and Koltun~\cite{sener2018moo} frame this as multi-objective optimization;
PCGrad~\cite{yu2020pcgrad} defines conflict as $\cos(\nabla\mathcal{L}_i, \nabla\mathcal{L}_j) < 0$
and projects conflicting gradients onto the normal plane;
CAGrad~\cite{liu2021cagrad} provides convergence guarantees under conflict.
In the DA literature specifically,
Du et al.~\cite{du2024gh} analyze obtuse-angle gradients between alignment and classification
losses in standard UDA and propose gradient rotation (GH/GH++),
while Phan et al.~\cite{phan2024pga} align per-domain gradients for prompt-based adaptation
of frozen vision-language models.
Our contribution extends gradient conflict analysis to the frozen-model DA setting.
We show that the cosine similarity
$\cos(\nabla_\theta\mathcal{L}_\text{task}, \nabla_\theta\mathcal{L}_\text{fid})$
between task and fidelity gradients---measured at the translator parameters $\theta$---is
a diagnostic that predicts method--task suitability:
$+0.84$ for mortality (aligned; delta translation suffices) versus $-0.21$ for sepsis
(conflicting; delta translation fails catastrophically).
Crucially, where PCGrad and CAGrad intervene \emph{algorithmically} by manipulating gradient
directions, our retrieval translator resolves the conflict \emph{architecturally}: the
cross-attention pathway over target-domain windows provides an alternative information
channel that reduces the dependence of the task gradient on input-space fidelity,
substantially mitigating the destructive interference without requiring gradient surgery
(detailed empirical gradient measurements in the full paper).

\medskip
\noindent
To summarize, no prior work combines a frozen target predictor, learned neural input-space
translation, instance-level target-domain retrieval, and gradient alignment theory for
clinical time-series domain adaptation.
The frozen-model constraint is motivated by real-world regulatory requirements, yet it
renders all existing clinical DA methods, general time-series DA methods, and foundation
model approaches inapplicable without fundamental modification.
Our retrieval-augmented translator fills this gap: it operates entirely in input space,
grounds translations in target-domain examples, and provides a theoretical framework
that explains when and why the approach succeeds.


% ============================================================
% Part 2: Problem Formulation + Method
% ============================================================
%% Part 2: Problem Formulation + Method
%% Sections 3–4 of the NeurIPS 2026 positioning paper

% ============================================================
\section{Problem Formulation}
\label{sec:formulation}
% ============================================================

\paragraph{Frozen-Model Translation Setting.}
Let $\mathcal{X}_S \sim P_S$ denote source-domain clinical time series (eICU) and $\mathcal{X}_T \sim P_T$ denote the target domain (MIMIC-IV).
A target-domain predictor $f_T : \mathbb{R}^{T \times F} \to \mathbb{R}^{T \times C}$ has been trained on $\mathcal{X}_T$ and is \emph{entirely frozen}: all parameters satisfy $\nabla_\phi f_T = 0$ and are never updated.
We learn a translator $g_\theta : \mathbb{R}^{T \times F} \to \mathbb{R}^{T \times F}$ that maps source inputs into the target input space so that the composite system $f_T \circ g_\theta$ performs well on source labels $y_S$.
The composite training objective is:
\begin{equation}
\label{eq:composite}
\mathcal{L}(\theta) \;=\; \mathcal{L}_{\mathrm{task}}\!\bigl(f_T(g_\theta(x_S)),\, y_S\bigr) \;+\; \lambda_{\mathrm{fid}}\,\mathcal{L}_{\mathrm{fid}}\!\bigl(g_\theta(x_S),\, x_S\bigr) \;+\; \lambda_{\mathrm{range}}\,\mathcal{L}_{\mathrm{range}}\!\bigl(g_\theta(x_S)\bigr),
\end{equation}
where $\mathcal{L}_{\mathrm{task}}$ is cross-entropy (classification) or MSE (regression) evaluated through the frozen predictor, $\mathcal{L}_{\mathrm{fid}}$ is an MSE fidelity penalty that regularises the translator toward identity, and $\mathcal{L}_{\mathrm{range}}$ penalises outputs outside the physiological range observed in $\mathcal{X}_T$.

The frozen constraint is motivated by regulatory and deployment considerations.
Under the FDA Software-as-a-Medical-Device (SaMD) ``locked algorithm'' pathway, the downstream predictor is a validated, auditable artifact whose parameters must not change post-deployment~\cite{fdaaiml2021}.
By learning only the input mapping $g_\theta$, clinical sites can adapt a certified predictor to local data distributions without re-validation of $f_T$.
This setting is strictly more constrained than test-time adaptation methods that update batch-normalisation statistics~\cite{wang2021tent} (which modify running mean and variance, altering the predictor's output for identical inputs) or partial model parameters~\cite{liang2020shot} (which retrain the feature extractor while freezing only the classifier head), and distinct from adapter-based approaches~\cite{houlsby2019adapters} that insert learnable modules into the predictor's computation graph. In our setting, the predictor is treated as a fixed, deterministic function---no internal parameter, statistic, or architectural element is modified.

\paragraph{Gradient Alignment Condition.}
A central difficulty of frozen-model translation is that the task gradient must flow \emph{through} the frozen predictor before reaching $\theta$.
Decomposing the gradient of Eq.~\ref{eq:composite}:
\begin{equation}
\label{eq:grad_decomp}
\nabla_\theta \mathcal{L} \;=\; \nabla_\theta \mathcal{L}_{\mathrm{task}} \;+\; \lambda_{\mathrm{fid}}\,\nabla_\theta \mathcal{L}_{\mathrm{fid}},
\end{equation}
we define the \emph{gradient alignment coefficient}:
\begin{equation}
\label{eq:alpha}
\alpha(\theta) \;=\; \cos\!\bigl(\nabla_\theta \mathcal{L}_{\mathrm{task}},\;\nabla_\theta \mathcal{L}_{\mathrm{fid}}\bigr).
\end{equation}

\begin{proposition}[Gradient Alignment Condition, informal]
\label{prop:alignment}
Stable frozen-model translation requires $\mathbb{E}[\alpha(\theta)] \geq 0$ during training.
When $\alpha < 0$, the fidelity and task gradients exhibit destructive interference: reducing $\mathcal{L}_{\mathrm{fid}}$ (moving the translation closer to identity) increases $\mathcal{L}_{\mathrm{task}}$, while reducing $\mathcal{L}_{\mathrm{task}}$ (improving predictions) increases domain divergence---violating the $\mathcal{H}\Delta\mathcal{H}$-divergence bound of \citet{bendavid2010}.
\end{proposition}

We verify this empirically across five clinical tasks.
Table~\ref{tab:alpha} reports gradient diagnostics measured at the translator parameters during the first epoch of training.\footnote{The full paper will extend these diagnostics across all three translator architectures and additional hyperparameter configurations, including per-epoch trajectories of $\alpha(\theta)$.}

\begin{table}[h]
\centering
\small
\caption{Gradient alignment diagnostics across tasks. $\|\nabla_\theta \mathcal{L}_{\mathrm{fid}}\| / \|\nabla_\theta \mathcal{L}_{\mathrm{task}}\|$ is the fidelity-to-task gradient magnitude ratio; $\alpha$ is the cosine alignment (Eq.~\ref{eq:alpha}). Positive label rate is per-timestep for AKI/Sepsis and per-stay for Mortality.}
\label{tab:alpha}
\begin{tabular}{lcccl}
\toprule
\textbf{Task} & \textbf{Label rate} & $\|\nabla \mathcal{L}_{\mathrm{fid}}\| / \|\nabla \mathcal{L}_{\mathrm{task}}\|$ & $\alpha$ & \textbf{Outcome} \\
\midrule
Mortality & 5.5\% (per-stay) & 2.8$\times$ & $+0.84$ & $+0.048$ AUROC \\
AKI & 11.95\% (per-ts) & moderate & $>0$ & $+0.056$ AUROC \\
Sepsis & 1.13\% (per-ts) & 5.7$\times$ & $-0.21$ & $+0.006$ AUROC (delta) \\
\bottomrule
\end{tabular}
\end{table}

Across the three classification tasks, the alignment coefficient is the strongest predictor of translation success: mortality ($\alpha = +0.84$) and AKI ($\alpha > 0$) succeed because fidelity acts as a step-size regulariser along the task gradient direction, whereas sepsis ($\alpha = -0.21$) suffers destructive interference.
Removing fidelity entirely ($\lambda_{\mathrm{fid}} = 0$) causes catastrophic divergence ($-0.101$ AUROC, from $0.719$ to $0.618$), confirming that the task gradient alone is too noisy to guide the translator---the fidelity term is necessary but, when misaligned, actively harmful.
This tension motivates the progression from delta-based translation (Section~\ref{sec:delta}) through shared-latent alignment (Section~\ref{sec:sl}) to retrieval-guided translation (Section~\ref{sec:retrieval}), where the retrieval mechanism provides an alternative information pathway that sidesteps the gradient alignment problem entirely.

% ============================================================
\section{Method}
\label{sec:method}
% ============================================================

We present three translator architectures of increasing capability. Each addresses a specific limitation of its predecessor, producing a clear ablation path: delta (limited by gradient alignment) $\to$ shared latent (limited by reconstruction bottleneck) $\to$ retrieval-guided (resolves both).

\paragraph{4.1\quad Delta Translator.}
\label{sec:delta}
The delta translator produces additive corrections:
\begin{equation}
\label{eq:delta}
g_\theta^{\delta}(x) \;=\; x + \Delta_\theta(x),
\end{equation}
where $\Delta_\theta$ is parameterised by a stack of \emph{AxialBlocks}---factored attention that alternates variable-wise attention (each timestep's $F$ features attend to one another) and temporal attention (each feature's $T$ timesteps attend across time).
This factorisation reduces the attention cost from $\mathcal{O}((TF)^2)$ to $\mathcal{O}(T \cdot F^2 + F \cdot T^2)$ while preserving full expressivity~\cite{ho2019axial}.
Static demographics (age, sex, height, weight) condition each block via FiLM layers~\cite{perez2018film}: $h \mapsto \gamma \odot h + \beta$ with $(\gamma, \beta)$ predicted from the static features.
Temporal attention is \emph{causal} (masked to prevent future information leakage) for per-timestep tasks (AKI, Sepsis, LoS) and \emph{bidirectional} for per-stay tasks (Mortality, KidneyFunction).

The residual formulation ensures $g_\theta$ begins near identity---early training produces small perturbations, avoiding the divergence observed with $\lambda_{\mathrm{fid}} = 0$.
However, the delta translator is fundamentally limited by Proposition~\ref{prop:alignment}: when $\alpha < 0$ (sepsis), the task and fidelity gradients create destructive interference that caps performance at $+0.006$ AUROC despite extensive hyperparameter search (13 approaches $\times$ 2 tasks; detailed in the full paper).

\paragraph{4.2\quad Shared Latent Translator.}
\label{sec:sl}
To address gradient misalignment, the shared latent (SL) translator replaces the direct input-to-input mapping with an encoder--decoder architecture that passes through a shared bottleneck $z \in \mathbb{R}^{T \times d_{\mathrm{latent}}}$.
Training proceeds in two phases.
\textbf{Phase~1} (autoencoder pretrain): the encoder--decoder is trained on MIMIC data only, with reconstruction and optional label prediction losses.
\textbf{Phase~2} (translation): the full loss combines task, reconstruction, distributional alignment, and label prediction:
\begin{equation}
\label{eq:sl_loss}
\mathcal{L}_{\mathrm{SL}} = \mathcal{L}_{\mathrm{task}} + \lambda_{\mathrm{recon}}\,\mathcal{L}_{\mathrm{recon}} + \lambda_{\mathrm{align}}\,\mathrm{MMD}(z_S, z_T) + \lambda_{\mathrm{label}}\,\mathcal{L}_{\mathrm{label\_pred}},
\end{equation}
where $\mathrm{MMD}$ is the maximum mean discrepancy with a multi-scale Gaussian kernel~\cite{gretton2012kernel} and $\mathcal{L}_{\mathrm{label\_pred}}$ provides a direct label supervision signal that bypasses the frozen predictor.
The MMD term provides distributional alignment at the latent level, replacing the input-space fidelity whose gradient opposes the task signal.

SL succeeds on tasks with sufficient label density: $+0.048$ AUROC ($+0.055$ AUCPR) on mortality (5.5\% per-stay labels), $+0.054$ AUROC ($+0.133$ AUCPR) on AKI (11.95\% per-timestep labels).
However, it fails on sepsis ($-0.017$ to $-0.043$ AUROC), exposing a \emph{reconstruction bottleneck}: the encoder--decoder must reconstruct all $T \times F$ features (e.g., $169 \times 48$) through the latent space.
At sepsis's 1.13\% positive rate, each training sequence contains ${\sim}1$ informative task gradient versus ${\sim}8{,}000$ reconstruction gradients---reconstruction dominates and the task signal cannot steer the latent representations.
Crucially, this is not a per-timestep or causality issue: AKI is also per-timestep and causal, yet SL succeeds because its 11.95\% label density provides $10\times$ more task gradient per sequence.
This controlled comparison isolates \emph{label density} as the root cause and motivates an architecture that does not rely on dense reconstruction to produce accurate translations.


\paragraph{4.3\quad Retrieval-Guided Translator.}
\label{sec:retrieval}
The retrieval-guided translator is our primary contribution.
Instead of global distributional alignment (SL's MMD) or direct input perturbation (delta), it performs \emph{instance-level retrieval}: each source timestep finds its nearest target-domain neighbours and selectively borrows feature-level information via cross-attention.
This opens a different information pathway: the model need not ``invent'' correct target-domain feature values from noisy task gradients, but can instead attend to actual target-domain exemplars.

\vspace{0.3em}
\noindent\textit{Memory Bank.}\quad
All target-domain (MIMIC) training data is encoded through the shared encoder and segmented into temporal windows of size $W$ (default $W{=}6$) with configurable stride $s$ for overlapping coverage.
Each window is mean-pooled to produce a compact descriptor:
\begin{equation}
\label{eq:memory}
M_j \;=\; \frac{1}{|W_j|} \sum_{t \in W_j} \mathrm{Enc}(x_T^{(i)})_t \;\in\; \mathbb{R}^{d_{\mathrm{latent}}},
\end{equation}
yielding a memory bank $M \in \mathbb{R}^{N_{\mathrm{windows}} \times d_{\mathrm{latent}}}$ that resides on GPU for fast retrieval.
The memory bank is rebuilt every $R$ epochs (default $R{=}3$), since the encoder evolves during training and a stale bank degrades retrieval quality.
Overlapping windows (stride $s{=}3$, yielding ${\sim}2\times$ more windows) improve performance on AKI ($+0.003$ AUROC, $+0.010$ AUCPR), but excessive overlap (stride $s{=}1$, ${\sim}2.4$M windows) degrades results ($-0.002$ AUROC), suggesting a coverage--redundancy tradeoff where moderate overlap improves recall of relevant exemplars but excessive redundancy dilutes the nearest-neighbour signal.

\vspace{0.3em}
\noindent\textit{Per-Timestep k-NN Retrieval.}\quad
For each source timestep $t$, we form a backward-looking query by mean-pooling source latents from $[\max(0, t{-}W{+}1),\; t{+}1)$, ensuring causal consistency.
Learned importance weights $\sigma(\mathbf{w}) \in [0,1]^{d_{\mathrm{latent}}}$ (sigmoid-gated, per latent dimension) modulate the distance metric:
\begin{equation}
\label{eq:retrieval_dist}
d(q_t, M_j) \;=\; \bigl\|\sqrt{\sigma(\mathbf{w})} \odot (q_t - M_j)\bigr\|_2,
\end{equation}
and the $K$ nearest windows per timestep are retrieved (default $K{=}16$).
Full timestep-level latents from each retrieved window are gathered to form the context tensor $C \in \mathbb{R}^{B \times T \times KW \times d_{\mathrm{latent}}}$.
Retrieval is per-timestep rather than per-sequence: we find that local temporal matching (window $W{=}6$) significantly outperforms wider windows ($W{=}12$: $+0.046$ vs.\ $+0.056$ AUROC on AKI), consistent with clinical time series where local physiological states are more informative than global trajectory similarity.

\vspace{0.3em}
\noindent\textit{Cross-Attention Fusion.}\quad
Retrieved context is integrated through $L$ stacked cross-attention blocks.
Each block performs three stages:
\begin{equation}
\label{eq:cross_attn}
\begin{aligned}
h_t' &= \mathrm{CrossAttn}(Q{=}h_t,\; K{=}C_t,\; V{=}C_t) + h_t, \\
\tilde{h} &= \mathrm{CausalSelfAttn}(h') + h', \\
\hat{h} &= \mathrm{FFN}(\tilde{h}) + \tilde{h},
\end{aligned}
\end{equation}
where $h_t$ is the source latent at timestep $t$, $C_t$ contains the $KW$ retrieved target latents for that timestep, and all sub-layers use pre-LayerNorm residual connections.
Stage~1 (per-timestep cross-attention) performs selective, feature-level borrowing from target exemplars---critically, this is \emph{learned}, not interpolated.
Compared to the $k$NN-MT approach of \citet{khandelwal2021knnmt}, which interpolates uniformly across all dimensions, cross-attention allows the model to attend selectively to different features from different neighbours.
Stage~2 (global causal self-attention) enforces temporal coherence across the translated sequence.

The number of cross-attention blocks $L$ is task-dependent and interacts with label density.
On AKI (11.95\% positive rate), increasing from $L{=}2$ to $L{=}3$ yields $+0.009$ AUROC---the additional capacity integrates richer target context from the dense supervision signal.
On sepsis (1.13\% positive rate), $L{=}3$ \emph{hurts} by $-0.006$ AUROC: extra parameters overfit when gradient signal is sparse.
This interaction provides a clean ablation of how architectural capacity must be matched to task supervision density.

\vspace{0.3em}
\noindent\textit{Why Retrieval Resolves Gradient Alignment.}\quad
The fundamental limitation of delta and SL translators is that the task gradient must specify both \emph{which features} to correct and \emph{what values} to produce---all by backpropagating through the frozen model's opaque hidden states.
Retrieval decouples these two problems by introducing an orthogonal information channel.
The memory bank supplies actual target-domain feature values; the model need only learn a soft \emph{selection} over them via cross-attention weights.
This constrains the output space to weighted combinations of real target exemplars, yielding a better-conditioned optimisation problem: the task gradient specifies \emph{which} neighbours are useful (a selection problem over $K$ candidates), while the neighbours themselves supply the \emph{content} (correct feature values from the target distribution).
Because the fidelity objective is now computed against outputs that are already grounded in the target domain, its gradient no longer opposes the task gradient---both objectives benefit from selecting more relevant target exemplars.

A critical design choice is detaching the source latent before querying the memory bank: $q_t = \mathrm{Enc}(x_S)_t.\mathrm{detach}()$.
Since $k$-NN retrieval involves a non-differentiable $\operatorname{argmin}$, gradients cannot propagate through the retrieval step; detaching prevents spurious gradient signals from the bank lookup.
The encoder is still updated through the downstream cross-attention and decoding losses.

The empirical results confirm this analysis.
On sepsis, where $\alpha = -0.21$ renders delta translation ineffective ($+0.006$ AUROC) and SL collapses ($-0.017$ to $-0.043$), retrieval achieves $+0.051$ AUROC---compared to $+0.006$ for the best delta configuration and negative results for the shared latent translator, establishing retrieval as the only architecture that produces clinically meaningful improvement on this task.
On AKI, retrieval ($+0.056$ AUROC, $+0.161$ AUCPR) surpasses both delta and SL, reaching an absolute AUROC of $91.14$ and AUCPR of $72.9$---above even the in-domain MIMIC-native LSTM ($89.7$ AUROC, $66.5$ AUCPR)~\cite{vandewater2024yaib}.
The AUCPR improvement ($+16.1$ points) is nearly $3\times$ larger than the AUROC gain ($+5.6$), indicating that the translator disproportionately improves precision at high recall---the regime most relevant for clinical deployment where false alarms carry real cost.

\vspace{0.3em}
\noindent\textit{FeatureGate.}\quad
We introduce a \emph{FeatureGate} module, applicable across all three paradigms: learnable per-feature sigmoid weights $\sigma(w_f) \in [0,1]$ that modulate the reconstruction loss:
\begin{equation}
\label{eq:feature_gate}
\mathcal{L}_{\mathrm{recon}}^{\mathrm{gated}} = \sum_{f=1}^{F} \sigma(w_f) \cdot \bigl(x_{\mathrm{out}}^f - x_{\mathrm{target}}^f\bigr)^2.
\end{equation}
This automatically down-weights reconstruction error on features irrelevant to the downstream task (e.g., features the frozen model ignores), focusing capacity on task-critical dimensions.
FeatureGate provides consistent gains: on mortality, SL+FG achieves $+0.048$ AUROC ($+0.055$ AUCPR) vs.\ plain SL $+0.044$ AUROC ($+0.048$ AUCPR).

\paragraph{4.4\quad Training Procedure.}
\label{sec:training}
Training follows a two-phase protocol common to SL and retrieval translators.
\textbf{Phase~1} (autoencoder pretrain, 15 epochs on MIMIC): the encoder--decoder is trained with reconstruction and label prediction losses.
For retrieval, we introduce \emph{self-retrieval pretraining}: during Phase~1, a memory bank
is built from the encoder's own MIMIC representations and refreshed periodically.
This provides non-degenerate context to cross-attention blocks from the start of training;
without self-retrieval, cross-attention receives zero tensors during Phase~1, collapsing
to a trivial pass-through that wastes the pretraining phase for the cross-attention parameters.
\textbf{Phase~2} (retrieval-guided translation, 50 epochs): the full loss is:
\begin{equation}
\label{eq:full_loss}
\mathcal{L} = \mathcal{L}_{\mathrm{task}} + \lambda_{\mathrm{recon}}\,\mathcal{L}_{\mathrm{recon}} + \lambda_{\mathrm{range}}\,\mathcal{L}_{\mathrm{range}} + \lambda_{\mathrm{align}}\,\mathrm{MMD} + \lambda_{\mathrm{target}}\,\mathcal{L}_{\mathrm{task}}^{\mathrm{MIMIC}} + \lambda_{\mathrm{label}}\,\mathcal{L}_{\mathrm{label\_pred}},
\end{equation}
where $\mathcal{L}_{\mathrm{task}}^{\mathrm{MIMIC}}$ applies the same task loss to MIMIC data through the frozen predictor, preventing representational drift.
The memory bank is rebuilt every 3 epochs; learning rate follows cosine annealing~\cite{loshchilov2017sgdr}; gradient norms are clipped to $1.0$; early stopping with patience 15 monitors the primary metric.

All data follows the YAIB preprocessing pipeline~\cite{vandewater2024yaib}: standardized variable selection (48 physiological features common to both databases), forward-fill imputation with missing indicators, per-feature normalization, and consistent patient-level train/validation/test splits.
This ensures that distributional differences between eICU and MIMIC-IV reflect institutional variation rather than preprocessing artifacts.
On top of YAIB's per-database normalization, we apply \emph{cross-domain normalisation}: an affine transform that matches source per-feature mean and variance to target statistics prior to translation.
This is computed once from training data and stored in the checkpoint, ensuring consistent normalisation at inference.

Phase~1 checkpoints are reusable: because pretraining depends only on the target data and model architecture (not Phase~2 hyperparameters), a single pretrain checkpoint serves all downstream experiments for a given task.
We verify this empirically: across 9 AKI configurations varying $K$, $W$, $L$, and loss weights, the Phase~1 label prediction loss converges to $0.462$--$0.466$, confirming Phase~2 independence.


% ============================================================
% Part 3: Comparative Analysis and Contributions
% ============================================================
%% Part 3: Comparative Analysis and Contributions
%% NeurIPS 2026 positioning paper
%% Depends on: Part 1 (\ref{sec:related}), Part 2 (method/formulation sections)

\section{Comparative Analysis and Contributions}
\label{sec:comparison}

With the related work surveyed in Section~\ref{sec:related}, we now
position our work against prior art.
We first provide a structured comparison across key desiderata
(Section~\ref{sec:comparison_table}), then present focused head-to-head analyses
with each method category (Section~\ref{sec:head_to_head}), situate our
empirical results against reference baselines (Section~\ref{sec:results_position}),
enumerate our concrete contributions (Section~\ref{sec:contributions}),
and conclude with an honest assessment of the current scope and limitations
(Section~\ref{sec:limitations}).


%% ============================================================
\subsection{Systematic Comparison}
\label{sec:comparison_table}

Table~\ref{tab:comparison} compares representative methods across eight
dimensions that collectively define the desiderata for clinical time-series
domain adaptation with a frozen predictor.
No prior method satisfies all eight.

\begin{table}[t]
\centering
\caption{Systematic comparison of domain adaptation and frozen-model methods
across key desiderata for clinical time-series translation.
\textbf{Model Frozen}: whether the target predictor's weights are entirely
fixed (\cmark), partially fixed (P), or trainable (\xmark).
\textbf{Adapt. Space}: Input (I), Representation (R), Weight (W), or Output (O).
\textbf{Sample-Cond.}: whether the transformation depends on the input instance.
\textbf{Clinical Tasks}: number of distinct clinical prediction tasks evaluated.
}
\label{tab:comparison}
\small
\setlength{\tabcolsep}{3.2pt}
\begin{tabular}{@{}lcccccccc@{}}
\toprule
\textbf{Method} &
\rotatebox{60}{\textbf{Model Frozen?}} &
\rotatebox{60}{\textbf{Adapt. Space}} &
\rotatebox{60}{\textbf{Sample-Cond.?}} &
\rotatebox{60}{\textbf{Retrieval?}} &
\rotatebox{60}{\textbf{Var.-Len. TS?}} &
\rotatebox{60}{\textbf{Clinical Tasks}} &
\rotatebox{60}{\textbf{Grad. Theory?}} \\
\midrule
\multicolumn{8}{@{}l}{\textit{Clinical DA}} \\
VRADA~\cite{purushotham2017vrada}   & \xmark & R & \cmark & \xmark & \cmark & 1 & \xmark \\
OTTEHR~\cite{ottehr2025}            & \xmark & R & \xmark & \xmark & \xmark & 3 & \xmark \\
Mutnuri et al.~\cite{mutnuri2024}   & \xmark & W & \cmark & \xmark & \cmark & 3 & \xmark \\
\midrule
\multicolumn{8}{@{}l}{\textit{General TS-DA}} \\
CoDATS~\cite{wilson2020codats}       & \xmark & R & \cmark & \xmark & \xmark & --- & \xmark \\
RAINCOAT~\cite{he2023raincoat}       & \xmark & R & \cmark & \xmark & \xmark & --- & \xmark \\
CLUDA~\cite{ozyurt2023cluda}         & \xmark & R & \cmark & \xmark & \xmark & --- & \xmark \\
\midrule
\multicolumn{8}{@{}l}{\textit{Frozen-model / Input-space}} \\
VPT~\cite{jia2022vpt}                & \cmark & I & \xmark & \xmark & \xmark & 0 & \xmark \\
Time-LLM~\cite{jin2024timellm}       & \cmark & I & \xmark & \xmark & \xmark & 0 & \xmark \\
TATO~\cite{qiu2026tato}              & \cmark & I & P      & \xmark & \cmark & 0 & \xmark \\
\midrule
\multicolumn{8}{@{}l}{\textit{Retrieval-augmented}} \\
$k$NN-MT~\cite{khandelwal2021knnmt}  & \cmark & O & \cmark & \cmark & \cmark & 0 & \xmark \\
TS-RAG~\cite{ning2025tsrag}          & \xmark & R & \cmark & \cmark & \cmark & 0 & \xmark \\
\midrule
\multicolumn{8}{@{}l}{\textit{Foundation model}} \\
ICareFM~\cite{burger2025icarefm}     & \xmark & W & \cmark & \xmark & \cmark & 4 & \xmark \\
\midrule
\textbf{Ours (Retrieval Translator)} & \cmark & \textbf{I} & \cmark & \cmark & \cmark & \textbf{5} & \cmark \\
\bottomrule
\end{tabular}
\vspace{-6pt}
\end{table}

Several observations stand out.
Representation-space (R) and weight-space (W) methods are architecturally coupled to the model's internals---they cannot be formulated without access to intermediate representations or weight tensors, violating the frozen-model constraint.
In contrast, input-space translation treats the target model as a \emph{black-box oracle} $f: \mathcal{X} \to \mathcal{Y}$.
While our efficient implementation leverages differentiable access (backpropagation through the frozen model) for gradient-based optimization, the method requires only query access to the model's input--output mapping.
In principle, this makes input-space translation applicable even to proprietary models exposed only via API, using gradient-free optimization (e.g., evolutionary strategies or Bayesian optimization)---a property no R-space or W-space method can satisfy, though we note that gradient-free training of the translator remains untested and would likely require substantially more queries.
The frozen-model methods from NLP and vision (VPT, Time-LLM) also operate in input space but use input-independent perturbations, limiting their capacity for heterogeneous clinical data.
$k$NN-MT~\cite{khandelwal2021knnmt} is closest in spirit---frozen model, retrieval,
per-timestep conditioning---but operates in output space (interpolating token distributions)
rather than input space (transforming the data itself).
Our retrieval translator is the only method that combines a fully frozen predictor,
input-space adaptation, sample-conditional neural translation, target-domain retrieval,
variable-length time-series support, evaluation across five clinical tasks, and a
theoretical gradient alignment framework.


%% ============================================================
\subsection{Head-to-Head Analysis}
\label{sec:head_to_head}

\paragraph{Against clinical DA methods.}
VRADA~\cite{purushotham2017vrada} pioneered adversarial alignment for clinical
time series but requires end-to-end training and was evaluated on a single task
(mortality).
Mutnuri et al.~\cite{mutnuri2024} address three of our five tasks (mortality,
AKI, LoS) with inductive transfer and standard adversarial DA, but fine-tune the
target model---precisely the operation our setting forbids.
On the shared AKI task, our frozen-model retrieval translator achieves
AUROC 91.14 and AUCPR 72.9, surpassing the MIMIC-native LSTM baseline (89.7 AUROC, 66.5 AUCPR) that represents
the performance ceiling when the model is trained on in-domain data---despite never modifying any predictor weights.
OTTEHR~\cite{ottehr2025} applies optimal transport to the same eICU--MIMIC
databases but performs population-level feature alignment without temporal
modeling, treating each feature independently.
Our instance-level, per-timestep retrieval and causal attention preserve the
temporal structure essential for sequential clinical tasks (AKI, sepsis, LoS).
In summary, prior clinical DA methods either sacrifice model integrity
(fine-tuning) or temporal structure (population-level alignment); we sacrifice
neither.

\paragraph{Against adapted DA baselines (DANN/CORAL/CoDATS).}
A natural question is whether standard DA methods can be adapted to our
frozen-model setting by learning an input transformation $T_\theta(\mathbf{x})$
that precedes the frozen predictor, with DANN's gradient reversal~\cite{ganin2016dann}
or CORAL's covariance alignment~\cite{sun2016coral} applied to the predictor's
internal representations.
We implement these adapted baselines (detailed in the full paper) and
find that they achieve limited improvement.
The structural argument is that clinical domain shift is \emph{heterogeneous}:
across 48 physiological features, the distributional shifts differ in magnitude,
direction, and functional form (e.g., heart rate may require a simple affine
correction while lab values exhibit nonlinear, context-dependent shifts).
A shallow input transform with a single domain discriminator or covariance
penalty lacks the capacity for such feature-specific, sample-conditional
mappings.
Our translator architectures provide structured inductive biases---additive
residuals (delta), latent bottleneck (shared latent), or target-domain cross-attention
(retrieval)---that decompose this complex mapping into learnable components.
The gradient alignment analysis (Section~\ref{sec:formulation}) further
explains the failure mode: when $\cos(\nabla_\theta \mathcal{L}_\text{task},
\nabla_\theta \mathcal{L}_\text{fid}) < 0$, as observed for sparse-label tasks,
the domain-alignment and task objectives pull the input transform in
opposing directions, a conflict that generic input transforms cannot resolve.

\paragraph{Against frozen-model methods.}
Prompt tuning~\cite{lester2021prompt} and VPT~\cite{jia2022vpt} learn a fixed
set of $\mathcal{O}(1)$ parameters (soft tokens or visual prompts) that are
shared across all inputs.
Clinical EHR time series vary from 1 to 169 timesteps, span 48--292 features with
complex missingness patterns, and exhibit patient-specific distributional
characteristics---far exceeding the capacity of input-independent perturbations.
Time-LLM~\cite{jin2024timellm} reprograms frozen LLMs for time-series forecasting,
but performs cross-\emph{modality} adaptation (language$\to$time series) for a
generic pre-trained model, not cross-\emph{domain} adaptation for a task-specific
validated clinical predictor.
TATO~\cite{qiu2026tato}, a concurrent ICLR~2026 contribution, is the closest
work in spirit: it explicitly proposes ``adapt data to model'' for frozen
time-series foundation models.
However, TATO relies on handcrafted transformation operators (slicing,
normalization, outlier correction) rather than learned neural translation,
targets forecasting rather than clinical prediction, and provides no
theoretical analysis of when data adaptation succeeds or fails.
We view TATO as validating the broader ``adapt data, not model'' paradigm;
our work extends it with (i)~learned end-to-end translation,
(ii)~retrieval-augmented cross-attention, and (iii)~gradient alignment theory.

\paragraph{Against retrieval-augmented methods.}
$k$NN-MT~\cite{khandelwal2021knnmt} is our closest architectural ancestor:
both systems pair a frozen model with per-timestep $k$-NN retrieval from a
domain-specific datastore.
The critical distinction lies in \emph{how} retrieval is used.
$k$NN-MT interpolates \emph{output token distributions}---blending the frozen
decoder's next-token prediction with the empirical distribution from retrieved
neighbors.
Our retrieval translator uses retrieved target-domain embeddings as
\emph{cross-attention context} in a Transformer decoder that transforms the
\emph{input}, enabling selective feature-level borrowing rather than
output-level interpolation.
Moreover, $k$NN-MT operates within a single language (e.g., domain-adapted
German$\to$English), while we bridge fundamentally different data distributions
with heterogeneous recording practices and device calibrations.
TS-RAG~\cite{ning2025tsrag} and RATD~\cite{liu2024ratd} retrieve time-series
segments for forecasting augmentation within the same domain; our retrieval
bridges \emph{across domains}---a qualitatively different use of the retrieval
paradigm.
RAM-EHR~\cite{xu2024ramehr} retrieves textual medical knowledge from PubMed
and knowledge graphs; we retrieve encoded \emph{time-series windows}
representing what the frozen predictor expects to see, grounding the translation
in the target domain's statistical structure rather than external knowledge.


%% ============================================================
\subsection{Results Positioning}
\label{sec:results_position}

Table~\ref{tab:results_position} situates our best results against
reference baselines from the YAIB benchmark~\cite{vandewater2024yaib}.
Full experimental details and ablations will appear in the full paper.

\begin{table}[t]
\centering
\caption{Our best translator results vs.\ reference baselines from YAIB~\cite{vandewater2024yaib}.
Classification metrics: AUROC and AUCPR $\times 100$; regression: MAE in natural units.
\textbf{Bold} indicates our result matches or surpasses the reference.
$\dagger$: MIMIC-native LSTM surpassed.
All our results use a \emph{frozen} MIMIC-IV predictor with no weight updates.}
\label{tab:results_position}
\small
\setlength{\tabcolsep}{3.2pt}
\begin{tabular}{@{}lcccccccc@{}}
\toprule
& \multicolumn{2}{c}{\textbf{Mortality}} & \multicolumn{2}{c}{\textbf{AKI}} & \multicolumn{2}{c}{\textbf{Sepsis}} & \textbf{LoS} & \textbf{KF} \\
& AUROC & AUCPR & AUROC & AUCPR & AUROC & AUCPR & MAE (h) & MAE (mg/dL) \\
\midrule
Frozen baseline  & 80.79 & 50.2 & 85.58 & 56.8 & 71.59 & 3.0 & 42.5 & 0.403 \\
eICU-native LSTM & 85.5  & 35.7 & 90.2  & 69.9 & 74.0  & 4.0 & 39.2 & 0.28 \\
MIMIC-native LSTM& 86.7  & 41.0 & 89.7  & 66.5 & 82.0  & 8.0 & 40.6 & 0.28 \\
\midrule
\textbf{Ours (best)}
  & \textbf{85.55} & \textbf{55.7} & $\mathbf{91.14}^\dagger$ & $\mathbf{72.9}^\dagger$ & \textbf{76.71} & \textbf{5.2} & \textbf{39.2} & 0.382 \\
$\Delta$ vs.\ frozen
  & $+$4.76 & $+$5.5 & $+$5.56 & $+$16.1 & $+$5.12 & $+$2.2 & $-$3.3 & $-$0.021 \\
\bottomrule
\end{tabular}
\vspace{-4pt}
\end{table}

\paragraph{Where we exceed expectations.}
On AKI, our frozen-model translator achieves AUROC 91.14, surpassing not only
the eICU-native LSTM (90.2) but also the \emph{MIMIC-native LSTM} (89.7)---the
model trained on in-domain data with full weight optimization.
The AUCPR results are even more striking: $72.9$ vs.\ $66.5$ for the MIMIC-native LSTM and $69.9$ for the eICU-native LSTM, a $+16.1$ point improvement over the frozen baseline.
AUCPR is particularly informative for clinical tasks with class imbalance because it measures the precision--recall tradeoff directly: a higher AUCPR means the model identifies more true positives without generating excessive false alarms.
The fact that the AUCPR gain ($+16.1$) is nearly $3\times$ the AUROC gain ($+5.6$) indicates that translation disproportionately improves the model's ability to distinguish positive cases at clinically relevant operating points.
On mortality, translation pushes AUCPR to $55.7$---far above both the eICU-native LSTM ($35.7$) and the MIMIC-native LSTM ($41.0$).
Notably, the frozen baseline already exceeds both native models on mortality AUCPR ($50.2$), and the translator extends this advantage further.
This demonstrates that input-space translation can recover performance
\emph{beyond} what the target model achieves on its own training distribution,
suggesting that the translated source data may introduce beneficial distributional
diversity that complements the frozen predictor's training distribution.
All three classification tasks surpass the eICU-native LSTM on both AUROC and AUCPR, confirming that
the translator closes the domain gap rather than merely compensating for it.
On LoS (regression), the translator achieves MAE of 39.2~hours, matching the
eICU-native LSTM exactly and improving upon the MIMIC-native LSTM (40.6~hours).
Cross-server reproducibility is strong: across 3 seeds on 2 physically
distinct GPU servers, AKI AUROC varies by only $\pm$0.0005, confirming that
the gains are robust and not artifacts of hardware-specific numerical behavior.

\paragraph{Where we acknowledge gaps.}
We identify three areas where performance does not yet match the strongest
reference points, and in each case, we characterize the underlying cause.
\emph{Sepsis}: the retrieval translator achieves $+5.12$ AUROC and $+2.2$ AUCPR, surpassing the eICU-native LSTM on both metrics (76.71 vs.\ 74.0 AUROC; 5.2 vs.\ 4.0 AUCPR---a 30\% relative AUCPR improvement), but the gains are smaller in magnitude than on denser tasks.
With only 1.13\% positive labels per timestep, the task gradient is extremely sparse: our gradient alignment analysis ($\alpha = -0.21$ for sepsis vs.\ $+0.84$ for mortality) correctly predicts the relative difficulty and identifies label sparsity as the primary bottleneck for further improvement.
\emph{Kidney Function}: the KF task involves 292 features including 192
generated cumulative statistics (\texttt{\_min\_hist}, \texttt{\_max\_hist},
\texttt{\_mean\_hist}, \texttt{\_count}) that must maintain internal
consistency after translation---a substantially harder translation target.
Our MAE of 0.382 improves over the frozen baseline (0.403) but a gap to the
eICU-native reference (0.28) remains.
\emph{Single domain pair}: our current evaluation covers eICU$\to$MIMIC-IV;
extension to HiRID$\to$MIMIC-IV is ongoing and will test generalization of
the approach across a third institutional distribution.


%% ============================================================
\subsection{Contributions}
\label{sec:contributions}

Our work makes five concrete contributions:

\begin{enumerate}[leftmargin=*,itemsep=2pt]

\item \textbf{Frozen-model preservation as a first-class constraint.}
We formalize the setting where the \emph{entire} target predictor is frozen---not
just the classifier head (SHOT~\cite{liang2020shot}) or a subset of layers---and
show that meaningful domain adaptation is achievable purely through input-space
translation.
This is directly motivated by the FDA's Software as a Medical Device (SaMD)
framework~\cite{fda2021samd}, which provides a simpler regulatory pathway for
``locked'' algorithms whose weights are not modified post-deployment.

\item \textbf{Input-space modularity.}
The translator is an external, independently trainable module that can be
swapped, updated, or audited without touching the validated predictor.
This decoupling enables a deployment model where the predictor and translator
have separate validation lifecycles---a practical requirement in regulated
clinical environments.

\item \textbf{Instance-level retrieval-augmented cross-attention for clinical
time-series DA.}
To the best of our knowledge, we are the first to combine per-timestep $k$-NN
retrieval from a target-domain memory bank with cross-attention decoding for
domain adaptation of clinical time series.
This provides each translated timestep with contextual grounding in what the
frozen model expects to see, enabling the fine-grained, feature-specific
transformations that generic input transforms cannot achieve.

\item \textbf{Universal paradigm across five clinical tasks.}
We evaluate on the broadest task set in the clinical DA literature:
three classification tasks (Mortality24, AKI, Sepsis) spanning per-stay and
per-timestep label structures, and two regression tasks (Length of Stay, Kidney
Function) with continuous targets.
The retrieval translator achieves the best or tied-best results on all five,
with improvements on both AUROC and AUCPR for all classification tasks---including
surpassing the MIMIC-native LSTM on both metrics for AKI.
This establishes retrieval translation as a universal paradigm with task-adapted capacity
(specifically, the number of cross-attention layers: 3 for dense-label tasks
like AKI, 2 for sparse-label tasks like sepsis).

\item \textbf{Gradient alignment as a predictive diagnostic.}
We introduce the cosine similarity
$\cos(\nabla_\theta \mathcal{L}_\text{task}, \nabla_\theta \mathcal{L}_\text{fid})$
as a diagnostic for frozen-model translation, connecting to the gradient conflict
literature in multi-task learning~\cite{yu2020pcgrad,liu2021cagrad} and
UDA~\cite{du2024gh,phan2024pga}.
This metric correctly predicts method--task suitability a priori
($+0.84$ for mortality, where delta translation succeeds; $-0.21$ for sepsis,
where it fails and retrieval is required) and reveals that the retrieval
architecture resolves the conflict \emph{structurally} through the cross-attention
information pathway, rather than requiring algorithmic gradient surgery.

\end{enumerate}


%% ============================================================
\subsection{Scope, Applicability, and Limitations}
\label{sec:limitations}

\paragraph{When to use input-space translation.}
Our approach is most applicable when three conditions hold:
(i)~a validated target-domain predictor exists and must remain frozen
(regulatory, institutional, or practical constraints);
(ii)~labeled data from both source and target domains is available for
translator training; and
(iii)~the target task provides sufficient label signal for the task gradient to guide translation. In our experiments, translation gains scale with label density: AKI (11.95\%) and mortality (5.5\%) yield the largest improvements, while sepsis (1.13\%) shows smaller but still meaningful gains ($+5.12$ AUROC), consistent with the gradient alignment analysis.
Under these conditions, the retrieval translator consistently closes the
domain gap across diverse clinical prediction tasks.

\paragraph{When alternatives are preferable.}
If the target model can be retrained, standard domain adaptation or foundation
model fine-tuning can optimize the predictor's own parameters---a strictly more
expressive search space---and may outperform input-space translation, though
this is not guaranteed: our translator surpasses the in-domain MIMIC-native
LSTM on AKI, suggesting that the inductive bias of retrieval-augmented
translation can sometimes compensate for the constrained parameter space.
If no target-domain data is available, source-free methods
(e.g., MAPU~\cite{ragab2023mapu}) are one option; alternatively,
our delta translator can operate without a target-domain memory bank,
using only the frozen model's task gradient and fidelity regularisation
to learn input corrections---achieving up to $+2.9$ AUROC on AKI
(vs.\ $+5.6$ with retrieval) when gradient alignment is positive ($\alpha > 0$).
If real-time inference latency is critical, the memory bank $k$-NN lookup
introduces approximately $2$--$3\times$ overhead compared to the frozen model
alone, which may be unacceptable in time-sensitive clinical alerting systems.

\paragraph{Current limitations.}
We identify five limitations that define the scope of the present work and
inform future directions:
\begin{enumerate}[leftmargin=*,itemsep=1pt]
\item \textit{Single domain pair.}
Our evaluation covers eICU$\to$MIMIC-IV.
Extension to additional source domains (HiRID$\to$MIMIC-IV) is under way
and will test whether the approach generalizes across institutional
distributions with different data collection protocols.

\item \textit{Label density dependence.}
Translation quality correlates with label density:
AKI (11.95\% positive rate) shows the largest gains ($\Delta$AUROC $+$5.56, $\Delta$AUCPR $+$16.1),
while sepsis (1.13\%) shows smaller but still meaningful improvements ($+5.12$ AUROC, $+2.2$ AUCPR).
The gradient alignment diagnostic quantifies this relationship, but does
not resolve the underlying information bottleneck for extremely sparse tasks.

\item \textit{Calibration degradation.}
Translation improves discrimination (AUROC, AUCPR) but can degrade probability
calibration (ECE), requiring post-hoc temperature scaling---an additional
validation step in clinical deployment.

\item \textit{Inference overhead.}
The retrieval translator requires maintaining a GPU-resident memory bank and
performing $k$-NN search at each timestep, adding computational and memory
costs relative to direct model inference.

\item \textit{Architecture specificity.}
Our frozen baselines are LSTMs from the YAIB
framework~\cite{vandewater2024yaib}.
Since input-space translation treats the predictor as a black-box function
$f: \mathcal{X} \to \mathcal{Y}$, the paradigm is architecture-agnostic in
principle; empirical validation with other target architectures (Transformers,
GRUs, temporal convolutional networks) remains future work.
\end{enumerate}

\noindent
We view these as open problems rather than fundamental obstacles.
The frozen-model, input-space translation paradigm opens a new direction in
clinical DA research, and the gradient alignment framework provides guidance
for extending it to new domains, tasks, and architectures.


% ============================================================
% Bibliography
% ============================================================
{\small
\bibliography{references}
}

\end{document}
